%!TEX program = lualatex
\documentclass[ngerman,headheight=70pt]{scrartcl}
\usepackage[ngerman]{babel}
\usepackage{blindtext}
\usepackage[utf8]{luainputenc}
\usepackage{scrlayer-scrpage}
\usepackage[T1]{fontenc}
\usepackage{fontspec}
\usepackage{multicol}
\usepackage{enumitem}
\usepackage[german=quotes]{csquotes}
\usepackage[bottom=2cm,footskip=8mm]{geometry}
\setmainfont{Source Sans Pro}
\MakeOuterQuote{"}

\setlength{\parindent}{0pt}
\setlength{\parskip}{6pt}

\pagestyle{plain.scrheadings}
\thispagestyle{scrheadings}
\lohead{\vspace*{-1.5cm}\includegraphics[scale=0.4]{../images/stupa_logo_rot.pdf}}
\rohead{\\Präsidium des\\Studierendenparlaments}

\lofoot[Universität Hamburg · Präsidium des Studierendenparlaments\\\hrule Von-Melle-Park 5 · D-20146 Hamburg]{Universität Hamburg · Präsidium des Studierendenparlaments\\\hrule Von-Melle-Park 5 · D-20146 Hamburg}
\cofoot[]{}
\cohead[\thepage]{}

\renewcommand{\thesubsection}{TOP \arabic{subsection}}
\renewcommand{\thesubsubsection}{\arabic{subsubsection}.}
\setcounter{subsection}{-1}

\begin{document}
    UHH · StuPa-Präsidium · Von-Melle-Park 5 · D-20146 Hamburg

    %%% PAGE ONE (EDIT LATER) %%%%
    % Hier den Titel des Protokolls eintragen
    \section*{Protokoll der 2. Sitzung des Studierendenparlaments vom 12. Mai 2016}

    % Protokollanten eintragen
    \textbf{Protokoll: Jim Martens}\\
    \textbf{Ort: S08 in der ehemaligen HWP}\\
    \textbf{Beginn: 18.15 Uhr}\\
    \textbf{Ende:  Uhr}

    % Wohl die komplizierteste Stelle im Protokoll
    % Als Faustregel sollten nicht mehr als 3 Namen auf einer Zeile stehen
    % Einfach an dieser Vorlage orientieren
    % Im Idealfall ist hier nicht viel zu ändern, da kaum jemand fehlt.
    % Rücktritte müssen natürlich pro Protokoll angepasst werden
    \vspace{0.5cm}
    \begin{tabular}{ll}
        Anwesend: & \\
            RCDS (5 Sitze): & Ramon Weilinger, Antonia Niecke, Ramin Shakiba, \\
                            & Jennifer Maack, Benjamin Welling \\
             CampusGrün (14 Sitze): & Laura Franzen, Geoffrey Youett, Elena Rysikova, \\
                                   & Philipp Droll, Melf Johannsen,\\
                                   & Tahnee Herzig, Mario Moldenhauer, Jim Martens,\\
                                   & Svenja Horn, Mirzo Ulugbek Khatamov,\\
                                   & Martin Sievert \\
             Bier-Liste (2 Sitze) : & Jakob Pape, Karen Martirosian \\
             WiWi (2 Sitze): & Claas-Friso Hente, Kay Zöllmer \\
             Unicorns (5 Sitze): & Katharina Kucza, Johannes Peplow, Annkathrin Löffler, \\
                                 & Andreas Hartkamp, Marielle Hermstrüwer \\
             Liste LINKS (3 Sitze): & Gunhild Berdal, Till Petersen, Sinah Mielich \\
             HWP (2 Sitze): & Samet Gunay, Ajdina Karahasan \\
             MIN (4 Sitze): & Ailina Salten, Thea Wahlers, Jan Detampel \\
             SDS* (3 Sitze): & Mena Winkler, Jacob Petersein, Artur Brückmann \\
             Bart-LISTE (2 Sitze): & Timo Zeimet, Dominic Laumer \\
             LHG (1 Sitz): & Tobias Heisig \\
             harte zeiten (1 Sitz): & Tobias Berking \\
             Jura (1 Sitz): & \\
             AL (2 Sitze): & Karima Schulze, Henri Weber \\
            & \\
        Entschuldigt: & Freya Schmitz (CampusGrün), Yasemin Günther (CG), Armin Günther, \\
                      & Lotte Rullkötter (MIN) \\
                                &\\
        Unentschuldigt abwesend: & Johann Baumhoefener (Jura) \\
                                &\\
        Rücktritte: & \\
    \end{tabular}
    \newpage
    %%% Vorgeschlagene Tagesordnung %%%
    \underline{Vorgeschlagene Tagesordnung}
    \begin{enumerate}[label={\textbf{Top \theenumi}},leftmargin=*]
        % folgende Zeile bei Sitzungsfortsetzungen einkommentieren und x durch
        % die Zahl des zuletzt behandelten TOPs ersetzen
        %\setcounter{enumi}{x}
        \item Alle Frauen*-Wahl
            \begin{enumerate}
                \item Bestätigung der Wahlniederschrift (V1617-020)
                \item Bestätigung der Referent*innen (V1617-019)
            \end{enumerate}
        \item Gründung des FSR kath. Theologie (V1617-017) (10 Min.)
        \item Antrag auf VV (V1617-018) (30 Min.)
        \item EM 2016 (V1617-014) (30 Min.)
        \item Verlegung der StuPa-Wahl in das Sommersemester (V1617-015) (30 Min.)
        \item Umgang mit der AfD (V1617-016) (30 Min.)
        \item Verschiedenes
    \end{enumerate}

    \newpage

    % Hier beginnt Teil A. Bei Fortsetzungen kommt hier nur der Teil über
    % die Feststellung der endgültigen Fassung der Tagesordnung.
    % Für sonstige Sitzungen ist dies bereits das Gerüst für die Formalia.
    % Man muss nur noch die protokollierten Inhalte eintragen.

    {\Large \underline{Teil A}}

    \subsection{Formalia}

    \subsubsection{Geschäftsbericht Präsidium}

    \subsubsection{Anfragen an das Präsidium}

    \subsubsection{Geschäftsbericht AStA}

    \subsubsection{Anfragen an den AStA}

    \subsubsection{Dringlichkeitsanträge des AStA}

    \subsubsection{Aktuelle Stunde (falls entsprechender Antrag vorliegt)}

    \subsubsection{Feststellung der endgültigen Fassung des Teils B der Tagesordnung}

    %%% BESCHLOSSENE TAGESORDNUNG %%%
    \underline{Endgültige Fassung der Tagesordnung}
    \begin{enumerate}[label={\textbf{Top \theenumi}},leftmargin=*]
        % folgende Zeile bei Sitzungsfortsetzungen einkommentieren und x durch
        % die Zahl des zuletzt behandelten TOPs ersetzen
        %\setcounter{enumi}{x}
        \item Alle Frauen*-Wahl
            \begin{enumerate}
                \item Bestätigung der Wahlniederschrift (V1617-020)
                \item Bestätigung der Referent*innen (V1617-019)
            \end{enumerate}
        \item Gründung des FSR kath. Theologie (V1617-017) (10 Min.)
        \item Antrag auf VV (V1617-018) (30 Min.)
        \item EM 2016 (V1617-014) (30 Min.)
        \item Verlegung der StuPa-Wahl in das Sommersemester (V1617-015) (30 Min.)
        \item Umgang mit der AfD (V1617-016) (30 Min.)
        \item Verschiedenes
    \end{enumerate}

    \subsubsection{Feststellung der Beschlussfähigkeit}

    \subsubsection{Genehmigung der Protokolle der vorangegangenen Sitzungen}

    \vspace{1cm}
    {\Large \underline{Teil B}}

    % Teil B beginnt
    % Dieser wird sich für jede Sitzung unterscheiden.
    % Jeder TOP ist eine subsection.
    % Vorhandene Beispiele anfangs entfernen und durch tatsächliche TOPs ersetzen

    % Folgende Zeile bei Sitzungsfortsetzungen einkommentieren und x mit
    % der Zahl des zuletzt behandelten TOPs ersetzen.
    %\setcounter{subsection}{x}
    \subsection{Alle Frauen*-Wahl}

    a) Bestätigung der Wahlniederschrift

    b) Bestätigung der Referent*Innen

    \subsection{Gründung des FSR kath. Theologie}

    \subsection{Antrag auf VV}

    \subsection{EM 2016}

    \subsection{Verlegung der StuPa-Wahl in das Sommersemester}

    \subsection{Umgang mit der AfD}

    \subsection{Verschiedenes}

\end{document}
